\begin{abstract}
  Some interruptions of a tool path are chosen by the planner; others are forced by the slicer's structure. We describe a slicer as a chain of maps, from a source through the part and the machine to the built part, and use the chain to locate forced interruptions and design around them. The obstruction is topological: a plan that must visit $c$ path components is interrupted at least $c-1$ times, and the chain identifies the link that supplies them. At the source, every reviewed solid-filling slicer traverses separate layers or surfaces, so its stops grow with resolution. Replacing the source by one connected volume, traversed by a three-dimensional Hilbert curve and carried into the part by optimal transport, gives a continuous path with the prescribed occupation measure on convex parts. In a computational prototype, it makes no stops on convex parts, where a layered raster's stops roughly double with resolution; on non-convex parts, a source split into near-convex pieces trades a fixed number of transitions for bounded excursions. At the machine, the components are the posture sheets of the robot's safe configurations, and a slicer that fixes the tool path first inherits the posture switches that path forces. Reversing the construction order gives kinematics-first slicing, in which the sheets guide the source split. On a three-link planar arm among obstacles, it makes the single posture reset required by the bound in a sampled availability model, while a geometry-first baseline makes two. The two constructions serve as demonstrations of the framework: one counting argument bounds both stops and posture switches, and in both a source split matched to the components attains the minimum.
\end{abstract}
